\subsection*{三组输入分别怎样演算}

先在串行、记录完整保留、一次只处理一个请求的条件下推演。每组历史都从账户 1 余额 10000 分、业务记录为空开始，\code{paper-001} 成交价为 3000 分。服务受理首个组合身份时，保存已确认输入与对应购买；后续同身份请求先核对输入。

第一组，R1 携带 K1 和 \code{paper-001}。服务提交 P1、O1、D1，保存成功结果：扣款 3000，扣后余额 7000。响应丢失，R2 继续携带 K1 和相同输入。服务找到 K1 的原记录，输入一致，返回 P1、O1、D1 以及当次扣后余额 7000，账户仍为 7000。返回的 7000 属于 P1 的结果；若稍后另有购买，当前余额需要另行查询。

第二组，R1 仍然以 K1 购买 \code{paper-001}，提交后余额为 7000。客户端误把 K1 配给另一商品 \code{paper-002}。服务找到 K1，发现决定效果的输入不同，返回身份输入冲突，P1 仍然存在，余额仍为 7000。它应让客户端定位原操作并修正本地状态。悄悄返回 \code{paper-001} 的成功结果，会令客户端误以为新商品已经买到；按新参数再执行，则会让同一身份表示两个效果。

第三组，Alice 明确另买一份 \code{paper-001}，客户端生成并保存 K2。请求体与第一次完全相同，组合身份不同。服务受理 P2、O2、D2，再扣 3000，余额从 7000 变为 4000。两笔购买都能单独查询。这个结果正是两次合法购买所需要的行为。

\begin{figure}[htbp]
\centering
\begin{tikzpicture}[x=1cm,y=1cm]
\node[booknote,text width=12cm] at (6.2,.9) {三组各自从余额 10000、空记录开始，第一笔均以 K1 成交 3000};
\node[booknode,text width=3.5cm] (a) at (1.7,-.3) {第二次：K1\\仍买 paper-001};
\node[booknode,text width=3.5cm] (b) at (6.2,-.3) {第二次：K1\\改买 paper-002};
\node[booknode,text width=3.5cm] (c) at (10.7,-.3) {第二次：K2\\仍买 paper-001};
\node[booknode,text width=3.5cm] (ar) at (1.7,-2) {返回原 P1、O1、D1\\余额保持 7000};
\node[bookresult,text width=3.5cm] (br) at (6.2,-2) {同身份输入冲突\\余额保持 7000};
\node[booknode,text width=3.5cm] (cr) at (10.7,-2) {建立 P2、O2、D2\\余额变为 4000};
\draw[bookarrow] (a)--(ar);
\draw[bookalert] (b)--(br);
\draw[bookarrow] (c)--(cr);
\node[booknote,text width=12cm] at (6.2,-3.35)
 {同键同参恢复原购买；同键异参保护原绑定；新键同参允许另一笔合法购买};
\end{tikzpicture}
\caption{身份与输入需要一起判断，三组情况产生不同处理结果。}
\label{fig:ch03-three-inputs}
\end{figure}


这三组规则把两类错误分开。身份遗失使应当汇合的尝试分成两项购买，造成漏识别；误复用使两项业务意图进入同一身份，造成错误合并。输入比较能发现其中一部分误复用。两次合法购买若恰好输入完全相同，服务无法区分客户端有意重试与误复用旧键，因此仍需客户端遵守新操作使用新键的规则。

可以把服务保存的绑定写为
\[
 M(S)=(V,X,P,R(S)).
\]
\(S\) 是完整身份；\(V\) 是输入解释规则的版本；\(X\) 是按该规则得到的规范化输入；\(P\) 是购买编号；\(R(S)\) 是已经保存的原业务结果。\(M\) 表示从身份查找完整绑定的映射，\(R(S)\) 只表示其中的结果部分。没有受理过的身份尚无对应项，成功完成或明确拒绝以后才具有相应终态。映射一旦形成，恢复同一 \(S\) 就应核对原 \(V,X\)，返回原 \(P,R(S)\)；新输入不能改写这项绑定。

这项规则可以用一对例子检查。K1、\code{paper-001} 对应 P1 以后，K1、\code{paper-002} 要报告冲突；K2、\code{paper-001} 仍可以建立另一项绑定。由此，身份决定查找哪项操作，输入决定该请求是否符合已建立的绑定。把输入摘要直接当作操作身份，会把两项输入相同的合法购买合并。

\begin{figure}[htbp]
\centering
\begin{tikzpicture}[x=1cm,y=1cm]
\node[booknode,text width=3.7cm] (k) at (1.8,0) {完整身份 S1\\账户、动作、K1};
\node[booknode,text width=6.8cm] (m) at (8,0) {已保存绑定 $M(S_1)$\\版本 V1、输入 X1、购买 P1、原结果 $R(S_1)$};
\draw[bookarrow] (k)--(m);
\node[booknode,text width=3.7cm] (same) at (1.8,-1.9) {恢复 S1、X1\\命中原绑定};
\node[bookresult,text width=4.2cm] (diff) at (7,-1.9) {恢复 S1、X2\\输入不同，报告冲突};
\node[booknode,text width=3.5cm] (new) at (11.6,-1.9) {新身份 S2、X1\\可另建购买};
\draw[bookarrow] (same.north)--(m.south west);
\draw[bookalert] (diff.north)--(m.south);
\node[booknote,text width=12cm] at (6.5,-3.3)
 {已保存绑定保持稳定；输入指纹用于核对绑定，输入相同的不同身份仍可各自成立};
\end{tikzpicture}
\caption{身份决定定位哪项操作，输入核对决定本次请求是否符合原绑定。}
\label{fig:ch03-binding}
\end{figure}


\subsection*{比较哪些输入}

先看输入里还缺什么。客户端确认时，\code{paper-001} 售价 3000 分；第一次请求迟到，目录已经涨为 3500 分。只有商品编号，服务就缺少用户接受哪个价格的依据。从这里起，请求增加报价编号 \code{quote\_id}：服务提供 Q7，记录该商品、数量 1、CNY 3000、适用账户或定价范围与截止时间，客户端确认后保存 Q7。服务首次成交时验证这份报价，恢复已成交操作时则读取原快照。

直接比较原始 JSON 字节会把无关的编码差异当成购买变化。字段顺序不同、空白不同，或明确写 \code{quantity=1} 与省略默认数量，都可能表达相同输入。另一些看似微小的差异会改变效果，例如币种从 CNY 换为 USD（美元），商品标识变了大小写，而商品目录把大小写视为不同 ID。

服务应先完成语法解析、字段校验和业务规范化。规范化是按已公布的业务规则把多种等价写法转成同一份可比较数据，例如把省略的 \code{quantity} 补成整数 1；它还要保留会改变购买效果的差别。服务随后再得到参与比较的输入。

本章固定规则为：\code{product\_id} 按商品目录的精确标识解释；\code{quantity} 缺省为 1，明确提供时也只能为整数 1；\code{currency} 缺省为 CNY，明确提供时只能为 CNY；未知字段一律拒绝；重复 JSON 成员名一律拒绝。字符串类型的 \code{"1"} 不按整数 1 接受。这样可把本章输入表示为一组已校验字段，不依赖原报文排版。

字段顺序可以规范化，业务集合和业务序列则要分别规定。以后购物车含多项商品时，数组顺序可能只用于展示，也可能决定优惠分摊、优先级或配送顺序；只有业务明确视为无序的部分才适合排序。金额转成最小货币单位时，还需要币种小数位和舍入规则。本章直接采用整数分，服务拒绝浮点金额与不支持的币种，减少输入解释的分叉。

\begin{table}[htbp]
\centering\small
\begin{tabularx}{\linewidth}{@{}p{4.3cm}p{2.5cm}X@{}}
\toprule
输入变化 & 本章处理 & 理由\\
\midrule
JSON 字段调换顺序 & 视为相同 & 对象成员顺序不决定购买效果\\
省略数量与数量为整数 1 & 视为相同 & 采用已公布的默认值\\
数量为字符串 \code{"1"} & 校验失败 & 类型不满足整数要求\\
商品标识改变 & 同键冲突 & 购买商品改变\\
未知字段出现 & 校验失败 & 本版没有该字段的业务解释\\
同名字段出现两次 & 校验失败 & 防止解析层选取不同值\\
\bottomrule
\end{tabularx}
\caption{输入相等由接口的业务解释规则确定。}
\end{table}

规范化规则属于特定接口。以后增加优惠券、收货地址、接受价格上限或发票选项时，要判断它们是否影响本次操作的效果，把相关字段纳入比较。诊断请求 ID、客户端等待超时和 trace 字段通常不改变购买效果，可以留在指纹之外。字段是否进入指纹应由业务含义决定，通用 JSON 排序工具无法完成这项判断。

服务可以存规范化输入，也可以存它的摘要以加速比较。摘要是由输入计算出的固定长度值；这里把用于识别本次输入的摘要称为指纹。输入变化通常会得到不同指纹，服务据此快速发现同键请求的输入冲突。摘要指纹应包含算法标识和输入规范版本，例如 \code{purchase-input-v1}、\mbox{\code{SHA-256}}。前者标明采用第几版字段解释规则，后者是一种输出 256 位摘要的算法；规则版本变了，即使算法相同，生成指纹时比较的内容也可能改变。

保存可审计的规范化字段，有助于在冲突时指出变化的是商品还是报价；若只留摘要，则只能判断不一致，事后解释能力更弱。摘要本身依赖抗碰撞假设：我们假定找到两份不同输入、使它们得到同一摘要，在实际计算能力下十分困难。这项假设给工程比较提供依据，有限长度摘要仍无法为任意长输入提供数学上绝无碰撞的证明。

把规范化过程记为 \(C_V\)，把摘要计算记为 \(H\)，则输入 \(x\) 的指纹为 \(H(C_V(x))\)。\(V\) 指定缺省值、类型和字段解释；\(C_V\) 得到可比较的业务输入；\(H\) 把这份输入压成固定长度摘要。两份报文经同版规范化后相同，指纹也相同。即使这个摘要计算完全正确，两个不同购买意图使用同一份输入，仍会得到同一指纹。加密摘要解决不了前面缺失的意图信息。

\begin{figure}[htbp]
\centering
\begin{tikzpicture}[x=1cm,y=1cm]
\node[booknode,text width=4cm] (a) at (1.9,0) {请求 A\\paper-001、Q7\\省略数量与币种};
\node[booknode,text width=4cm] (b) at (1.9,-2.3) {请求 B\\paper-001、Q7\\数量 1、CNY};
\node[booknode,text width=4.2cm] (v) at (7,-1.15) {按同一版本规范化\\商品、整数 1、CNY、Q7\\得到相同字段组 X1};
\node[booknode,text width=2.9cm] (h) at (11.8,-1.15) {固定编码\\计算指纹 H(X1)};
\draw[bookarrow] (a)--(v);
\draw[bookarrow] (b)--(v);
\draw[bookarrow] (v)--(h);
\node[bookresult,text width=4cm] (c) at (1.9,-4.4) {请求 C 改为 Q8\\会改变报价依据};
\node[bookresult,text width=4.2cm] (cx) at (7,-4.4) {得到不同字段组 X2\\同身份应报告冲突};
\draw[bookalert] (c)--(cx);
\node[booknote,text width=12cm] at (6.6,-5.6)
 {解析时拒绝重复成员名；输入核对和实际执行共用同一份已校验数据};
\end{tikzpicture}
\caption{先确定哪些写法表达同一输入，再计算摘要；固定长度指纹保留算法与规则版本。}
\label{fig:ch03-normalization}
\end{figure}


部署 v2 后，旧 K1 仍应按保存的 \code{purchase-input-v1} 解释。若默认数量或字段含义已经改变，服务应调用旧版本规则，或明确拒绝不兼容恢复，保留查询原结果的路径。把旧摘要与新规则临时计算的摘要直接比较，会把版本变化混入用户输入变化。

可以具体看两份输入的转换。请求 A 只有 \code{product\_id=paper-001}、\code{quote\_id=Q7}。请求 B 另写 \code{quantity=1}、\code{currency=CNY}，并调整成员顺序。解析并校验后，两者都转换为同样的有序字段组：版本 \code{purchase-input-v1}，商品 \code{paper-001}，数量整数 1，币种 CNY，报价 Q7。指纹从这组定义明确的数据编码计算，因而一致。

请求 C 把报价改为 Q8，它的字段组在报价位置不同，服务据此报告冲突。编码还应明确字段边界与类型，例如采用带长度和类型的信息，防止把两个相邻字段直接拼接成相同字符串。对一份输入，规范化比较与实际执行还应共用同一份已校验数据。

重复成员名需要在解析时发现。如果请求写了两次 \code{quote\_id}，某一层选第一个值而另一层选最后一个值，就可能出现指纹绑定 Q7、执行却使用 Q8 的错配。网关、业务层和持久层各自重新解释未经校验的原报文，会使同一次处理内部出现两种购买输入。身份记录因此还要保存动作版本、指纹版本和必要的规范化输入。

\subsection*{价格变化时，原输入与原成交各是什么}

前面的请求只写商品，3000 分来自服务端商品目录。假设 R1 已按 3000 分成交，响应丢失，目录随后涨价到 3500 分。R2 携带 K1 到达，服务应先找到原记录，核对原请求含义，再返回 P1 原成交 3000 分的结果。把当前价格 3500 重新纳入同键请求的指纹，会把一次正常恢复误判为参数冲突。

对已完成购买，成交金额来自原记录；对尚未受理的新购买，计价规则决定何时确定金额。如果 R1 根本没有到达，R2 是服务第一次收到 K1，而且此时价格已变为 3500 分，那么仅有 \code{product\_id} 的请求没有留下用户接受 3000 分的服务端证据。接口若允许按受理时价格成交，K1 可能首次成交 3500。用户在界面看到的旧价格与服务承诺的计价时点必须一致。

前面加入的报价引用解决这个歧义。客户端把已经确认的 Q7 纳入本地操作记录和规范化输入。服务首次受理时核验报价是否可用，成交后保存 Q7 的价格快照。客户端请求体不直接决定扣款额，服务依据已经验证的报价确定金额。

\begin{figure}[htbp]
\centering
\begin{tikzpicture}[x=1cm,y=1cm]
\node[booknode,text width=3.5cm] (q) at (1.8,0) {Alice 确认 Q7\\paper-001、3000 分};
\node[booknode,text width=3.8cm] (p) at (6.6,0) {K1 在 Q7 有效期内成交\\P1 保存成交快照};
\node[booknode,text width=3.7cm] (r) at (11.5,0) {目录变为 3500\\恢复 K1 仍返回 3000};
\draw[bookarrow] (q)--(p);
\draw[bookarrow] (p)--(r);
\node[bookresult,text width=3.8cm] (late) at (6.6,-2)
 {K1 首次到达\\Q7 已过期\\保存报价过期结果};
\draw[bookalert] (q.south)--++(0,-2)--(late.west);
\node[booknode,text width=3.7cm] (new) at (11.5,-2) {Alice 接受 Q8 的 3500\\新确认、新身份 K2};
\draw[bookarrow] (late)--(new);
\node[booknote,text width=12cm] at (6.7,-3.4)
 {报价期限决定新操作能否成交；已完成操作按原成交结果恢复};
\end{tikzpicture}
\caption{原输入、首次受理条件和已提交结果分别对应不同判断时点。}
\label{fig:ch03-quote}
\end{figure}


如果 Q7 过期前 P1 已提交，过期后重试 K1，原成交仍为 3000 分；这次重试查的是已完成购买。若 K1 第一次到达时 Q7 已过期，则返回明确的报价过期结果，未产生订单和扣款。Alice 接受 Q8 的 3500 分后，应创建 K2，再提交 Q8。把 K1 的 \code{quote\_id} 从 Q7 改成 Q8 会触发同键异参冲突，因为报价已经进入用户确认内容。

这还区分了输入固定与结果固定。K1、Q7 固定了用户提交的购买依据，能否成交取决于首次有效受理时的业务条件；一旦结果提交，后续恢复引用该结果。报价过期会影响尚未受理的购买，已经提交的 P1 则继续按成交快照解释。把这两个时点混在一起，会让重试重新取得一个由当前环境计算的结果。

报价有效期与幂等恢复有效期也是两条时间规则。前者约束新购买能否按某价格成交，后者约束旧身份还能否用于恢复。若接口支持在首次受理时锁定报价，锁定记录本身也要可靠保存；否则进程重启后重新查价，仍可能失去第一次确认的金额。这里要求业务成交快照与结果一起保存。原子写入指这些记录作为一个整体生效：共同提交，或共同回滚。

\subsection*{拒绝结果与临时故障}

同一身份还需要说明失败后是否继续代表原操作。账户只有 2000 分，K1 按 Q7 支付 3000 分时，可以保存余额不足终态，并在结果中以 \code{rejected} 表达明确拒绝。终态表示这项操作已经作出最终业务结论，后续相同身份的恢复返回原结论；充值后的新购买则使用 K2，第 5.4 节会给出完整报文。

上述路径已经保存了明确拒绝，服务故障还可能使它来不及形成这样的记录。例如数据库暂时不可用、依赖服务超时或处理进程中断，可能发生在业务判断之前，也可能发生在提交结果的途中。服务可以返回暂时不可用或结果待确认，客户端保留 K1 后查询、等待或重试。若服务不知道事务是否提交，必须通过原身份恢复，不能把故障自动解释为可以安全创建新键。HTTP 5xx 指状态码处于 500 至 599 的服务器错误响应；它是否伴随业务生效，需要由实际提交与恢复路径确认。

\begin{figure}[htbp]
\centering
\begin{tikzpicture}[x=1cm,y=1cm]
\node[booknode,text width=3.5cm] (a) at (1.7,0) {语法校验失败\\明确尚未受理};
\node[booknode,text width=3.5cm] (b) at (6.2,0) {余额不足\\已提交拒绝终态};
\node[bookresult,text width=3.5cm] (c) at (10.7,0) {执行中失联\\尚未确认提交结果};
\node[booknode,text width=3.5cm] (ar) at (1.7,-1.9) {按规则修正输入\\再提交原身份};
\node[booknode,text width=3.5cm] (br) at (6.2,-1.9) {同身份返回原拒绝\\充值后新确认用新键};
\node[bookresult,text width=3.5cm] (cr) at (10.7,-1.9) {保留原身份与输入\\查询或重试查明事实};
\draw[bookarrow] (a)--(ar);
\draw[bookarrow] (b)--(br);
\draw[bookalert] (c)--(cr);
\end{tikzpicture}
\caption{恢复动作依赖受理与提交证据。一个泛化的失败提示不足以选择下一步。}
\label{fig:ch03-failure-evidence}
\end{figure}


参数语法错误发生在身份受理以前时，可以规定不占用 K1。客户端修正语法后仍使用 K1，前提是服务确实没有建立该身份的业务记录。若输入已被受理，服务后续发现余额不足或报价过期，并保存为终态，则该身份被绑定到原输入和结果。接口需要给出可辨别的错误码与是否已受理信息，客户端据此管理操作记录。

因此，相同的失败提示可能要求完全不同的恢复动作。明确未受理，允许按规则修正输入；已有终态，沿用原结论；提交尚不确定，保留原身份查明事实。决定下一步的依据，是系统留下了哪类持久记录，以及当前调用方能否读取它。
\FloatBarrier
